% Held out of the submitted version at the advisor's request (2026-09-08):
% restore with % Held out of the submitted version at the advisor's request (2026-09-08):
% restore with % Held out of the submitted version at the advisor's request (2026-09-08):
% restore with % Held out of the submitted version at the advisor's request (2026-09-08):
% restore with \input{sections/limitations_holdout} before the future-work paragraph in discussion.tex (Conclusion and Future Work)
\textbf{Limitations.}  The measurements come from one fleet of Jetson boards
on a wired LAN and from two models of one family, OPT-6.7B and OPT-350M.  Other model families, device classes, and lossy or intermittent links, under
which a missed KV column fails the contract into replay, remain to be
measured.  At 7B, Reconfigure was
measured at two positions.  The protection domain covers
fail-stop failures of non-head workers: the head stage and the coordinator
are unprotected, a last-stage failure falls back to replay because no
downstream survivor can confirm the final slot, and a failure during prefill
has no prefix to reconstruct.  The injected failure is a deterministic exception raised after the victim's
input for position $P$ has reached the coordinator.  Process kills,
out-of-memory terminations, failures before the mirrored input lands
(contract check (iv)), and the replay fallback they trigger were not
measured, so the reported latencies characterize the parity path rather than
the average over all failure timings.  Latency was measured with
synchronous forwarding; asynchronous forwarding adds a heartbeat-bounded
detection delay common to all families that we did not quantify, and a
worker that hangs rather than crashes can still be misattributed.  Backup weights are preloaded and Eq.~\ref{eq:mem} reserves memory for them;
a lazy variant that loads only the failed stage's weights at failure time
would relax the reservation at a recovery cost between \sys{} and
Reconfigure, a middle point we did not measure.  The coordinator holds all
parity state and is itself unprotected: its failure loses every in-flight
request's parity and mirror state, and a standby coordinator would need those
columns replicated to it, which we leave to future work.  Its memory grows
with concurrent requests and context length beyond the per-request figures
above.  A head failure falls to full-prefix replay on its preloaded backup, a
path we did not measure.  The
scheduler does not optimize where backups land, so the two-failure result
holds for mappings that, like the 7B one, place the two backups on distinct
nodes.  After a recovery the promoted backup serves unprotected until the
coordinator re-solves the mapping, a re-protection step we did not measure.
The flat-in-$P$ result is measured to $P=32$; the reconstruction loop is
linear in the prefix at a per-position cost the fit bounds below 12\,ms.  The fidelity probe covers
one tier pair and one two-layer stage on the 350M deployment.  The storage
advantage depends on placement geometry: double parity saves 46\% over
replication on the 7B placement but only 11\% on the head-heavy 350M
placement, where a third parity column would cost more state than
replication while tolerating fewer failures.  Each recovery-latency fit
pools three trials per position; the victim-sensitivity sweep, the 7B Reconfigure trials, and the 350M
sweeps are one trial per position; the protection-cost figures are means over
three rounds.  Network traffic and per-token energy were not measured.
Recovery latencies are measured with the failure detected on the request
RPC path; a failure detected only by heartbeat adds up to the 5\,s timeout.  Every
measurement is single-request; concurrency and prompt length were not
varied.  Recovery latency excludes post-recovery degradation: a promoted
backup that serves two stages decodes more slowly until re-protection.  The
$k=1$ and $k=2$ trials differ in victim identity and backup tier as well as
in $k$.

 before the future-work paragraph in discussion.tex (Conclusion and Future Work)
\textbf{Limitations.}  The measurements come from one fleet of Jetson boards
on a wired LAN and from two models of one family, OPT-6.7B and OPT-350M.  Other model families, device classes, and lossy or intermittent links, under
which a missed KV column fails the contract into replay, remain to be
measured.  At 7B, Reconfigure was
measured at two positions.  The protection domain covers
fail-stop failures of non-head workers: the head stage and the coordinator
are unprotected, a last-stage failure falls back to replay because no
downstream survivor can confirm the final slot, and a failure during prefill
has no prefix to reconstruct.  The injected failure is a deterministic exception raised after the victim's
input for position $P$ has reached the coordinator.  Process kills,
out-of-memory terminations, failures before the mirrored input lands
(contract check (iv)), and the replay fallback they trigger were not
measured, so the reported latencies characterize the parity path rather than
the average over all failure timings.  Latency was measured with
synchronous forwarding; asynchronous forwarding adds a heartbeat-bounded
detection delay common to all families that we did not quantify, and a
worker that hangs rather than crashes can still be misattributed.  Backup weights are preloaded and Eq.~\ref{eq:mem} reserves memory for them;
a lazy variant that loads only the failed stage's weights at failure time
would relax the reservation at a recovery cost between \sys{} and
Reconfigure, a middle point we did not measure.  The coordinator holds all
parity state and is itself unprotected: its failure loses every in-flight
request's parity and mirror state, and a standby coordinator would need those
columns replicated to it, which we leave to future work.  Its memory grows
with concurrent requests and context length beyond the per-request figures
above.  A head failure falls to full-prefix replay on its preloaded backup, a
path we did not measure.  The
scheduler does not optimize where backups land, so the two-failure result
holds for mappings that, like the 7B one, place the two backups on distinct
nodes.  After a recovery the promoted backup serves unprotected until the
coordinator re-solves the mapping, a re-protection step we did not measure.
The flat-in-$P$ result is measured to $P=32$; the reconstruction loop is
linear in the prefix at a per-position cost the fit bounds below 12\,ms.  The fidelity probe covers
one tier pair and one two-layer stage on the 350M deployment.  The storage
advantage depends on placement geometry: double parity saves 46\% over
replication on the 7B placement but only 11\% on the head-heavy 350M
placement, where a third parity column would cost more state than
replication while tolerating fewer failures.  Each recovery-latency fit
pools three trials per position; the victim-sensitivity sweep, the 7B Reconfigure trials, and the 350M
sweeps are one trial per position; the protection-cost figures are means over
three rounds.  Network traffic and per-token energy were not measured.
Recovery latencies are measured with the failure detected on the request
RPC path; a failure detected only by heartbeat adds up to the 5\,s timeout.  Every
measurement is single-request; concurrency and prompt length were not
varied.  Recovery latency excludes post-recovery degradation: a promoted
backup that serves two stages decodes more slowly until re-protection.  The
$k=1$ and $k=2$ trials differ in victim identity and backup tier as well as
in $k$.

 before the future-work paragraph in discussion.tex (Conclusion and Future Work)
\textbf{Limitations.}  The measurements come from one fleet of Jetson boards
on a wired LAN and from two models of one family, OPT-6.7B and OPT-350M.  Other model families, device classes, and lossy or intermittent links, under
which a missed KV column fails the contract into replay, remain to be
measured.  At 7B, Reconfigure was
measured at two positions.  The protection domain covers
fail-stop failures of non-head workers: the head stage and the coordinator
are unprotected, a last-stage failure falls back to replay because no
downstream survivor can confirm the final slot, and a failure during prefill
has no prefix to reconstruct.  The injected failure is a deterministic exception raised after the victim's
input for position $P$ has reached the coordinator.  Process kills,
out-of-memory terminations, failures before the mirrored input lands
(contract check (iv)), and the replay fallback they trigger were not
measured, so the reported latencies characterize the parity path rather than
the average over all failure timings.  Latency was measured with
synchronous forwarding; asynchronous forwarding adds a heartbeat-bounded
detection delay common to all families that we did not quantify, and a
worker that hangs rather than crashes can still be misattributed.  Backup weights are preloaded and Eq.~\ref{eq:mem} reserves memory for them;
a lazy variant that loads only the failed stage's weights at failure time
would relax the reservation at a recovery cost between \sys{} and
Reconfigure, a middle point we did not measure.  The coordinator holds all
parity state and is itself unprotected: its failure loses every in-flight
request's parity and mirror state, and a standby coordinator would need those
columns replicated to it, which we leave to future work.  Its memory grows
with concurrent requests and context length beyond the per-request figures
above.  A head failure falls to full-prefix replay on its preloaded backup, a
path we did not measure.  The
scheduler does not optimize where backups land, so the two-failure result
holds for mappings that, like the 7B one, place the two backups on distinct
nodes.  After a recovery the promoted backup serves unprotected until the
coordinator re-solves the mapping, a re-protection step we did not measure.
The flat-in-$P$ result is measured to $P=32$; the reconstruction loop is
linear in the prefix at a per-position cost the fit bounds below 12\,ms.  The fidelity probe covers
one tier pair and one two-layer stage on the 350M deployment.  The storage
advantage depends on placement geometry: double parity saves 46\% over
replication on the 7B placement but only 11\% on the head-heavy 350M
placement, where a third parity column would cost more state than
replication while tolerating fewer failures.  Each recovery-latency fit
pools three trials per position; the victim-sensitivity sweep, the 7B Reconfigure trials, and the 350M
sweeps are one trial per position; the protection-cost figures are means over
three rounds.  Network traffic and per-token energy were not measured.
Recovery latencies are measured with the failure detected on the request
RPC path; a failure detected only by heartbeat adds up to the 5\,s timeout.  Every
measurement is single-request; concurrency and prompt length were not
varied.  Recovery latency excludes post-recovery degradation: a promoted
backup that serves two stages decodes more slowly until re-protection.  The
$k=1$ and $k=2$ trials differ in victim identity and backup tier as well as
in $k$.

 before the future-work paragraph in discussion.tex (Conclusion and Future Work)
\textbf{Limitations.}  The measurements come from one fleet of Jetson boards
on a wired LAN and from two models of one family, OPT-6.7B and OPT-350M.  Other model families, device classes, and lossy or intermittent links, under
which a missed KV column fails the contract into replay, remain to be
measured.  At 7B, Reconfigure was
measured at two positions.  The protection domain covers
fail-stop failures of non-head workers: the head stage and the coordinator
are unprotected, a last-stage failure falls back to replay because no
downstream survivor can confirm the final slot, and a failure during prefill
has no prefix to reconstruct.  The injected failure is a deterministic exception raised after the victim's
input for position $P$ has reached the coordinator.  Process kills,
out-of-memory terminations, failures before the mirrored input lands
(contract check (iv)), and the replay fallback they trigger were not
measured, so the reported latencies characterize the parity path rather than
the average over all failure timings.  Latency was measured with
synchronous forwarding; asynchronous forwarding adds a heartbeat-bounded
detection delay common to all families that we did not quantify, and a
worker that hangs rather than crashes can still be misattributed.  Backup weights are preloaded and Eq.~\ref{eq:mem} reserves memory for them;
a lazy variant that loads only the failed stage's weights at failure time
would relax the reservation at a recovery cost between \sys{} and
Reconfigure, a middle point we did not measure.  The coordinator holds all
parity state and is itself unprotected: its failure loses every in-flight
request's parity and mirror state, and a standby coordinator would need those
columns replicated to it, which we leave to future work.  Its memory grows
with concurrent requests and context length beyond the per-request figures
above.  A head failure falls to full-prefix replay on its preloaded backup, a
path we did not measure.  The
scheduler does not optimize where backups land, so the two-failure result
holds for mappings that, like the 7B one, place the two backups on distinct
nodes.  After a recovery the promoted backup serves unprotected until the
coordinator re-solves the mapping, a re-protection step we did not measure.
The flat-in-$P$ result is measured to $P=32$; the reconstruction loop is
linear in the prefix at a per-position cost the fit bounds below 12\,ms.  The fidelity probe covers
one tier pair and one two-layer stage on the 350M deployment.  The storage
advantage depends on placement geometry: double parity saves 46\% over
replication on the 7B placement but only 11\% on the head-heavy 350M
placement, where a third parity column would cost more state than
replication while tolerating fewer failures.  Each recovery-latency fit
pools three trials per position; the victim-sensitivity sweep, the 7B Reconfigure trials, and the 350M
sweeps are one trial per position; the protection-cost figures are means over
three rounds.  Network traffic and per-token energy were not measured.
Recovery latencies are measured with the failure detected on the request
RPC path; a failure detected only by heartbeat adds up to the 5\,s timeout.  Every
measurement is single-request; concurrency and prompt length were not
varied.  Recovery latency excludes post-recovery degradation: a promoted
backup that serves two stages decodes more slowly until re-protection.  The
$k=1$ and $k=2$ trials differ in victim identity and backup tier as well as
in $k$.

